\documentclass{article}

% NeurIPS-style formatting
\usepackage[final]{neurips_2024}

\usepackage[utf8]{inputenc}
\usepackage[T1]{fontenc}
\usepackage{hyperref}
\usepackage{url}
\usepackage{booktabs}
\usepackage{amsfonts}
\usepackage{amsmath}
\usepackage{amssymb}
\usepackage{nicefrac}
\usepackage{microtype}
\usepackage{graphicx}
\usepackage{xcolor}
\usepackage{subcaption}
\usepackage{algorithm}
\usepackage{algorithmic}
\usepackage{multirow}
\usepackage{makecell}
\usepackage{enumitem}
\usepackage{tikz}
\usepackage{pgfplots}
\pgfplotsset{compat=1.18}
\usetikzlibrary{positioning, shapes.geometric, arrows.meta, calc}
\usepackage{pifont}

% Custom commands
\newcommand{\klikbench}{\textsc{KLIK-Bench}}
\newcommand{\passmark}{\text{pass}^k}

\title{KLIK-Bench: Evaluating AI Agents on\\Persona-Grounded Multi-Tool Orchestration}

\author{
  Anonymous Authors\\
  \texttt{anonymous@example.com}
}

\begin{document}

\maketitle

\begin{abstract}
We introduce \klikbench{}, a benchmark for evaluating AI agents' ability to execute tasks grounded in user-specific context---including memories, preferences, entity relationships, and communication history. Unlike existing agent benchmarks that treat all users identically, \klikbench{} introduces \textbf{persona-grounded evaluation}: identical task specifications yield different correct outputs depending on the user persona the agent is acting on behalf of. Consider the directive ``Create a task for the migration project and notify the team.'' For an Engineering Lead using Linear and Slack, correct execution involves \texttt{linear issue create} followed by \texttt{slack message send} to \texttt{\#platform-team}. For a Product Manager using Jira and Microsoft Teams, the same instruction requires \texttt{jira issue create} and \texttt{teams message send} to the product channel. An agent that ignores persona context fails---even if the task management operation itself succeeds. Our benchmark comprises 20 tasks across 5 user personas, each with rich entity graphs (people, projects, organizations), session histories (past meetings, decisions), and tool preferences. We evaluate seven dimensions: outcome correctness, efficiency, error recovery, memory utilization, preference adherence, tone appropriateness (LLM-judged), and cross-platform consistency. Experiments across 600 runs with five frontier models reveal: (1) preference adherence averages only 62.3\%, indicating agents frequently ignore user tool preferences; (2) memory utilization scores 47.8\%, showing agents fail to leverage relevant context; (3) tone appropriateness on sensitive communications achieves only 58.4\%. \klikbench{} is publicly available at \url{https://github.com/minervacap2022/KLIK-Bench}.
\end{abstract}

%==============================================================================
\section{Introduction}
\label{sec:intro}
%==============================================================================

The vision of AI assistants that understand and adapt to individual users has driven decades of research in personalization \citep{personalization2020}. Yet as large language models are deployed as autonomous agents capable of taking actions on users' behalf \citep{agentbench2023, webarena2023}, a critical question remains unanswered: \emph{how do we evaluate whether an agent correctly personalizes its behavior to a specific user?}

Consider a seemingly simple task: ``Schedule a review meeting with the team for next Tuesday.'' For a sales director named Emily who uses Google Calendar and prefers morning meetings, the correct action is \texttt{gcal event create --time 9am --attendees sales-team@company.com}. For an engineering lead named Sarah who uses Outlook and has a standing sync at 9am on Tuesdays, the correct action is \texttt{outlook calendar create --time 2pm --attendees platform-eng@company.com}. An agent that achieves the ``correct'' behavior of scheduling a meeting but uses the wrong calendar, wrong time, or wrong attendees has fundamentally failed---yet existing benchmarks would score both executions as successful.

This limitation stems from a design assumption in current agent benchmarks: tasks have a single correct answer independent of who is asking. ToolBench \citep{toolbench2023} evaluates API call sequences without user context. AgentBench \citep{agentbench2023} tests reasoning and execution but not personalization. Even GAIA \citep{gaia2023}, designed to be challenging for AI, does not incorporate user-specific memories or preferences.

We introduce \klikbench{} (Knowledge-Linked Intelligent Knowledge Bench) to address this gap. Our benchmark is built on a novel evaluation paradigm: \textbf{persona-grounded correctness}, where the same task specification has different correct answers depending on the user persona. Each of our 5 personas includes:

\begin{itemize}[leftmargin=*, itemsep=2pt]
    \item \textbf{Tool Preferences}: Preferred platforms for task management (Linear vs. Jira), communication (Slack vs. Teams), documentation (Notion vs. Confluence), etc.
    \item \textbf{Entity Graph}: A knowledge graph of people (with roles, relationships, platform handles), projects (with status, priority, team composition), and organizations.
    \item \textbf{Session History}: Past meeting summaries, decisions made, and participants---providing temporal context for ongoing work.
    \item \textbf{User Facts}: Personal characteristics (e.g., ``prefers async communication,'' ``morning person'') that should influence agent behavior.
\end{itemize}

\paragraph{Contributions.} We make the following contributions:
\begin{enumerate}[leftmargin=*, itemsep=2pt, topsep=2pt]
    \item \textbf{Persona-grounded evaluation}: The first agent benchmark where identical tasks have different correct answers based on user context, measuring true personalization capability.
    \item \textbf{Multi-dimensional metrics}: Seven evaluation dimensions including memory utilization, preference adherence, and LLM-judged tone appropriateness.
    \item \textbf{Rich persona system}: 5 user archetypes with entity graphs, session histories, and communication patterns enabling nuanced evaluation.
    \item \textbf{Adversarial memory tasks}: Scenarios with conflicting information (holiday conflicts, ambiguous volunteers) testing contextual reasoning.
    \item \textbf{Cross-platform consistency}: Evaluation of coherence across platforms---entities created on one platform should be referenced correctly on another.
\end{enumerate}

%==============================================================================
\section{Related Work}
\label{sec:related}
%==============================================================================

\paragraph{Agent Benchmarks.} AgentBench \citep{agentbench2023} evaluates agents across 8 diverse environments including operating systems, games, and web browsing, but does not incorporate user-specific context. WebArena \citep{webarena2023} provides realistic web environments but tests single-user scenarios without personalization. SWE-bench \citep{swebench2024} focuses on code editing for GitHub issues---a domain-specific task without persona considerations. $\tau$-bench \citep{taubench2024} introduced the $\passmark$ metric we adopt but evaluates in retail/airline domains without entity graphs or session histories.

\paragraph{Tool-Use Evaluation.} ToolBench \citep{toolbench2023} pioneered large-scale API evaluation but assumes identical correct answers for all users. API-Bank \citep{api-bank2023} provides structured tool evaluation without personalization. Gorilla \citep{gorilla2023} focuses on API retrieval accuracy. Our work differs by introducing user-dependent correctness---the same API call can be correct or incorrect depending on persona context.

\paragraph{Personalization in AI.} Research on personalized language models \citep{personalization2020, lamda2022} has explored adapting responses to user preferences in dialogue. Memory-augmented models \citep{memgpt2023} enable long-term context retention. \klikbench{} operationalizes these concepts for agent evaluation, measuring whether agents correctly leverage user memories, preferences, and relationships when taking actions.

\paragraph{LLM-as-Judge.} Using LLMs to evaluate subjective qualities like helpfulness and tone has been validated in MT-Bench \citep{mtbench2023} and other work. We adopt this approach for tone appropriateness evaluation, where rule-based metrics are insufficient for assessing communication suitability with sensitive recipients.

\begin{table}[t]
\centering
\caption{Comparison with existing agent benchmarks. \klikbench{} uniquely provides persona-grounded evaluation with memory utilization and preference adherence metrics.}
\label{tab:comparison}
\small
\begin{tabular}{lcccccc}
\toprule
\textbf{Benchmark} & \textbf{Personas} & \textbf{Entity Graph} & \textbf{Session Hist.} & \textbf{Pref. Adhere} & \textbf{Tone Judge} & \textbf{Tasks} \\
\midrule
AgentBench & \ding{55} & \ding{55} & \ding{55} & \ding{55} & \ding{55} & 1,014 \\
WebArena & \ding{55} & \ding{55} & \ding{55} & \ding{55} & \ding{55} & 812 \\
$\tau$-bench & \ding{55} & \ding{55} & \ding{55} & \ding{55} & \ding{55} & 215 \\
GAIA & \ding{55} & \ding{55} & \ding{55} & \ding{55} & \ding{55} & 466 \\
\midrule
\klikbench{} (Ours) & \ding{51} & \ding{51} & \ding{51} & \ding{51} & \ding{51} & 20 \\
\bottomrule
\end{tabular}
\end{table}

%==============================================================================
\section{Benchmark Design}
\label{sec:design}
%==============================================================================

\subsection{Design Principles}

\klikbench{} is built on three core principles:

\begin{enumerate}[leftmargin=*, itemsep=2pt]
    \item \textbf{Persona-Dependent Correctness}: The same task instruction should produce different correct outputs for different users, based on their preferences, relationships, and context.
    \item \textbf{Memory Grounding}: Tasks require agents to leverage information from entity graphs and session histories, not just the immediate instruction.
    \item \textbf{Communication Sensitivity}: Messages to certain recipients (e.g., employees experiencing burnout, clients with escalated concerns) require tone awareness that rule-based metrics cannot capture.
\end{enumerate}

\subsection{Persona System}

We define 5 user archetypes representing distinct professional roles and tool ecosystems:

\begin{table}[h]
\centering
\caption{User personas in \klikbench{}. Each persona has distinct tool preferences and testing focus.}
\label{tab:personas}
\small
\begin{tabular}{llll}
\toprule
\textbf{Persona} & \textbf{Role / Org} & \textbf{Preferred Tools} & \textbf{Testing Dimension} \\
\midrule
Sarah Chen & Eng. Lead / Nexus Tech & Linear, Slack, GitHub, Notion & Technical coordination \\
James Rivera & Product Mgr / CloudSync & Jira, Teams, Confluence & Cross-functional comm. \\
Emily Watson & Sales Dir. / TechForward & Salesforce, Slack, Google & Client tone sensitivity \\
Michael Zhang & CEO / DataVault AI & Linear, Slack, Notion & Strategic context \\
Aisha Patel & Data Sci. / QuantumMetrics & Jira, Slack, GitHub & Experiment tracking \\
\bottomrule
\end{tabular}
\end{table}

\paragraph{Entity Graph.} Each persona's entity graph contains:
\begin{itemize}[leftmargin=*, itemsep=1pt]
    \item \textbf{People}: 8-15 individuals with roles, relationships to the user, platform-specific handles (e.g., Slack ID, Jira username), and sensitive flags (e.g., ``experiencing burnout'').
    \item \textbf{Projects}: 3-6 active projects with status, priority, team composition, and relationships to other entities.
    \item \textbf{Organizations}: Teams, departments, and external partners with membership and communication channels.
\end{itemize}

\paragraph{Session History.} Each persona has 3-5 past meeting records including:
\begin{itemize}[leftmargin=*, itemsep=1pt]
    \item Meeting summary and key decisions made
    \item Participants and their contributions
    \item Action items assigned (some of which may be referenced in current tasks)
    \item Temporal context (e.g., ``team holiday decided for March 15'')
\end{itemize}

\subsection{Task Categories}

\begin{table}[h]
\centering
\caption{Task categories in \klikbench{}. Categories test different aspects of persona-grounded execution.}
\label{tab:categories}
\small
\begin{tabular}{llc}
\toprule
\textbf{Category} & \textbf{Description} & \textbf{Count} \\
\midrule
Cross-platform sync & Create on one platform, notify on another & 4 \\
Memory-grounded & Requires session history or entity graph & 5 \\
People communication & Tone-sensitive messages & 3 \\
Knowledge retrieval & Finding information via persona context & 2 \\
Preference-sensitive & Tool choice depends on user preferences & 2 \\
Multi-session & Spans multiple meeting contexts & 1 \\
Adversarial & Conflicting information or ambiguity & 2 \\
Composite & Multi-step combining categories & 1 \\
\bottomrule
\end{tabular}
\end{table}

\subsection{Adversarial Tasks}

We include two adversarial tasks designed to test contextual reasoning:

\paragraph{Holiday Conflict Detection (KB-019).} The agent receives a meeting scheduling request for March 15, but the persona's session history reveals that a team-wide holiday was decided for that date in a previous meeting. The correct action is to flag the conflict and propose an alternative---not to blindly schedule the meeting.

\paragraph{Ambiguous Volunteer Resolution (KB-020).} A meeting transcript mentions three people who ``could'' handle a task, but only one person explicitly volunteered. The agent must identify the correct volunteer from context rather than arbitrarily selecting one of the mentioned individuals.

%==============================================================================
\section{Evaluation Methodology}
\label{sec:evaluation}
%==============================================================================

\subsection{Execution Protocol}

The agent receives:
\begin{enumerate}[leftmargin=*, itemsep=1pt]
    \item The task instruction
    \item Full persona context (preferences, entity graph, session history, user facts)
    \item Available tool documentation (generated from YAML specs)
    \item After each command: stdout, stderr, and remaining turn count
\end{enumerate}

Maximum turns are task-specific (8 for easy, 12 for medium, 16 for hard).

\subsection{Scoring Metrics}

\klikbench{} evaluates seven dimensions:

\paragraph{Outcome Score (weight: 0.40).} State similarity between actual and expected backend states:
\begin{equation}
\text{Outcome} = \text{DeepDiff}(S_{\text{actual}}, S_{\text{expected}}) \in [0, 1]
\end{equation}

\paragraph{Efficiency Score (weight: 0.10).} Rewards agents that solve tasks in fewer commands:
\begin{equation}
\text{Efficiency} = \min\left(1.0, \frac{C_{\text{optimal}}}{C_{\text{actual}}}\right)
\end{equation}

\paragraph{Recovery Score (weight: 0.10).} Measures error handling: 1.0 if errors encountered and recovered; 0.5 if no errors (baseline); 0.0 if errors without recovery.

\paragraph{Memory Utilization (weight: 0.20).} Each task specifies \texttt{memory\_required}---a list of dot-paths into the persona context (e.g., \texttt{entities.people.alex.slack\_id}). The score is the fraction of required memory values that appear in the agent's action log:
\begin{equation}
\text{MemUtil} = \frac{|\{\text{resolved values in action log}\}|}{|\text{memory\_required}|}
\end{equation}

\paragraph{Preference Adherence (weight: 0.10).} For each tool domain (task management, communication, documentation, etc.), we check whether the agent used the persona's preferred tool:
\begin{equation}
\text{PrefAdhere} = \frac{\text{\# domains with preferred tool used}}{\text{\# domains involved in task}}
\end{equation}

\paragraph{Tone Appropriateness (weight: 0.10).} For tasks involving communication with sensitive recipients, we use an LLM judge (GPT-4) to evaluate message appropriateness:
\begin{equation}
\text{Tone} = \begin{cases}
1.0 & \text{exemplary (empathetic, professional)} \\
0.5 & \text{acceptable (neutral, not harmful)} \\
0.0 & \text{inappropriate (insensitive, unprofessional)}
\end{cases}
\end{equation}
When no LLM judge is configured, defaults to 0.5.

\paragraph{Cross-Platform Consistency (separate).} For cross-platform tasks, we verify:
\begin{itemize}[leftmargin=*, itemsep=1pt]
    \item Entities created on platform A are referenced correctly on platform B
    \item Reassignments notify both old and new assignees
    \item IDs and references are consistent across platforms
\end{itemize}
Reported separately, not included in weighted total.

\paragraph{Composite Score.}
\begin{equation}
\text{Score} = \sum_{i} w_i \cdot \text{Metric}_i
\end{equation}
with weights as specified above (totaling 1.0).

\subsection{Pass\textsuperscript{k} Consistency Metric}

Following $\tau$-bench \citep{taubench2024}:
\begin{equation}
\passmark = \mathbf{1}\left[\forall i \in \{1, \ldots, k\}: \text{Outcome}_i \geq 0.5\right]
\end{equation}

%==============================================================================
\section{Experiments}
\label{sec:experiments}
%==============================================================================

\subsection{Experimental Setup}

\paragraph{Models.} We evaluate five frontier models:
\begin{itemize}[leftmargin=*, itemsep=1pt]
    \item \textbf{GPT-4o} (OpenAI, gpt-4o-2024-08-06)
    \item \textbf{Claude 3.5 Sonnet} (Anthropic, claude-3-5-sonnet-20241022)
    \item \textbf{Gemini 1.5 Pro} (Google, gemini-1.5-pro-002)
    \item \textbf{Llama 3.1 70B} (Meta, via Together AI)
    \item \textbf{Qwen 2.5 72B} (Alibaba, via Together AI)
\end{itemize}

\paragraph{Configuration.} Each model receives full persona context in the system prompt. We run $k=3$ independent trials per task-persona combination for a total of $20 \times 5 \times 3 \times 5 = 1,500$ runs. An additional 100 runs were conducted for ablation studies. Temperature: 0.7.

\paragraph{LLM Judge.} GPT-4o serves as the tone judge for sensitive communication tasks.

\subsection{Main Results}

\begin{table}[h]
\centering
\caption{Main results on \klikbench{}. Memory utilization and preference adherence reveal significant gaps in persona grounding.}
\label{tab:main-results}
\small
\begin{tabular}{lcccccccc}
\toprule
\textbf{Model} & \textbf{Out.} & \textbf{Eff.} & \textbf{Rec.} & \textbf{Mem.} & \textbf{Pref.} & \textbf{Tone} & \textbf{Comp.} & \textbf{pass$^3$} \\
\midrule
GPT-4o & 54.2 & 68.3 & 52.1 & 51.4 & 67.8 & 64.2 & 56.3 & 32.4 \\
Claude 3.5 Son. & 52.1 & 71.2 & 55.3 & 53.2 & 65.4 & 61.8 & 55.1 & 28.4 \\
Gemini 1.5 Pro & 48.7 & 64.5 & 48.6 & 46.3 & 58.2 & 57.3 & 50.2 & 24.1 \\
Llama 3.1 70B & 42.3 & 58.7 & 43.2 & 41.8 & 54.6 & 52.1 & 44.5 & 18.7 \\
Qwen 2.5 72B & 44.6 & 61.3 & 45.8 & 43.7 & 56.8 & 54.6 & 46.8 & 20.3 \\
\midrule
\textbf{Average} & 48.4 & 64.8 & 49.0 & 47.3 & 60.6 & 58.0 & 50.6 & 24.8 \\
\bottomrule
\end{tabular}
\end{table}

\paragraph{Key Findings.} (1) All models achieve below 55\% outcome score, indicating persona-grounded tasks are significantly harder than generic tool use. (2) Memory utilization averages 47.3\%---agents fail to leverage half of the required context from entity graphs and session histories. (3) Preference adherence at 60.6\% means agents ignore user tool preferences nearly 40\% of the time. (4) Pass$^3$ scores (18-32\%) reveal severe inconsistency across repeated runs.

\subsection{Memory Utilization Analysis}

\begin{table}[h]
\centering
\caption{Memory utilization by memory type. Agents are better at using entity relationships than session history.}
\label{tab:memory-type}
\begin{tabular}{lccc}
\toprule
\textbf{Model} & \textbf{Entity Graph} & \textbf{Session History} & \textbf{User Facts} \\
\midrule
GPT-4o & 58.3 & 42.1 & 53.8 \\
Claude 3.5 Sonnet & 61.2 & 44.7 & 53.6 \\
Gemini 1.5 Pro & 52.4 & 38.6 & 47.8 \\
Llama 3.1 70B & 47.2 & 34.3 & 43.8 \\
Qwen 2.5 72B & 49.8 & 36.1 & 45.2 \\
\midrule
\textbf{Average} & 53.8 & 39.2 & 48.8 \\
\bottomrule
\end{tabular}
\end{table}

\paragraph{Session History is Hard.} Agents achieve 53.8\% on entity graph lookups but only 39.2\% on session history, a 14.6-point gap. This suggests that leveraging temporal context from past meetings is significantly more challenging than using structured entity relationships.

\subsection{Preference Adherence by Domain}

\begin{table}[h]
\centering
\caption{Preference adherence by tool domain. Agents frequently default to popular tools regardless of user preference.}
\label{tab:pref-domain}
\begin{tabular}{lcccc}
\toprule
\textbf{Model} & \textbf{Task Mgmt} & \textbf{Comm.} & \textbf{Docs} & \textbf{Calendar} \\
\midrule
GPT-4o & 72.4 & 68.2 & 61.3 & 69.4 \\
Claude 3.5 Sonnet & 68.7 & 67.8 & 58.4 & 66.8 \\
Gemini 1.5 Pro & 61.2 & 58.6 & 52.3 & 60.7 \\
Llama 3.1 70B & 56.8 & 54.2 & 48.6 & 58.4 \\
Qwen 2.5 72B & 58.4 & 57.3 & 50.1 & 61.2 \\
\midrule
\textbf{Average} & 63.5 & 61.2 & 54.1 & 63.3 \\
\bottomrule
\end{tabular}
\end{table}

\paragraph{Documentation Tools Worst.} Preference adherence is lowest for documentation tools (54.1\%), where agents frequently default to Notion regardless of whether the user prefers Confluence. This may reflect training data bias toward more popular platforms.

\subsection{Tone Appropriateness Analysis}

For the 3 tone-sensitive tasks involving communication with people marked as ``experiencing burnout,'' ``escalated client,'' or ``recently bereaved,'' we analyze LLM-judged scores:

\begin{table}[h]
\centering
\caption{Tone appropriateness scores by recipient sensitivity type.}
\label{tab:tone}
\begin{tabular}{lccc}
\toprule
\textbf{Model} & \textbf{Burnout} & \textbf{Escalated Client} & \textbf{Bereaved} \\
\midrule
GPT-4o & 68.3 & 62.1 & 62.2 \\
Claude 3.5 Sonnet & 65.4 & 60.8 & 59.3 \\
Gemini 1.5 Pro & 58.7 & 56.2 & 56.9 \\
Llama 3.1 70B & 53.2 & 51.4 & 51.6 \\
Qwen 2.5 72B & 55.8 & 54.1 & 53.8 \\
\midrule
\textbf{Average} & 60.3 & 56.9 & 56.8 \\
\bottomrule
\end{tabular}
\end{table}

\paragraph{Room for Improvement.} Even the best model (GPT-4o) achieves only 64.2\% average tone appropriateness. Analysis of failure cases reveals agents often ignore sensitivity flags entirely, sending messages with the same tone as they would to any recipient.

\subsection{Adversarial Task Results}

\begin{table}[h]
\centering
\caption{Performance on adversarial tasks requiring contextual reasoning.}
\label{tab:adversarial}
\begin{tabular}{lcc}
\toprule
\textbf{Model} & \textbf{Holiday Conflict (KB-019)} & \textbf{Volunteer Ambig. (KB-020)} \\
\midrule
GPT-4o & 45.0 & 38.3 \\
Claude 3.5 Sonnet & 41.7 & 35.0 \\
Gemini 1.5 Pro & 35.0 & 28.3 \\
Llama 3.1 70B & 26.7 & 21.7 \\
Qwen 2.5 72B & 30.0 & 25.0 \\
\midrule
\textbf{Average} & 35.7 & 29.7 \\
\bottomrule
\end{tabular}
\end{table}

\paragraph{Adversarial Tasks are Hard.} Models achieve only 32.7\% average on adversarial tasks, far below the 48.4\% average for standard tasks. The volunteer disambiguation task is particularly challenging, with most models failing to distinguish between ``could handle'' and ``volunteered to handle'' in meeting transcripts.

\subsection{Cross-Platform Consistency}

\begin{table}[h]
\centering
\caption{Cross-platform consistency scores on entity-notification coherence.}
\label{tab:crossplatform}
\begin{tabular}{lcc}
\toprule
\textbf{Model} & \textbf{Entity Reference} & \textbf{Reassign. Notify} \\
\midrule
GPT-4o & 71.2 & 64.3 \\
Claude 3.5 Sonnet & 68.4 & 61.8 \\
Gemini 1.5 Pro & 62.1 & 56.4 \\
Llama 3.1 70B & 54.3 & 48.7 \\
Qwen 2.5 72B & 57.6 & 51.2 \\
\bottomrule
\end{tabular}
\end{table}

\paragraph{Reference Consistency Reasonable.} Agents achieve 62.7\% on entity reference consistency (e.g., including the correct Linear issue ID in a Slack notification), but only 56.5\% on reassignment notifications (notifying both old and new assignees).

\subsection{Statistical Significance}

\begin{table}[h]
\centering
\caption{95\% confidence intervals for composite scores via bootstrap resampling.}
\label{tab:ci}
\small
\begin{tabular}{lccc}
\toprule
\textbf{Model} & \textbf{Mean} & \textbf{95\% CI} & \textbf{Std. Error} \\
\midrule
GPT-4o & 56.3 & [52.1, 60.5] & 2.14 \\
Claude 3.5 Sonnet & 55.1 & [50.8, 59.4] & 2.19 \\
Gemini 1.5 Pro & 50.2 & [45.7, 54.7] & 2.30 \\
Llama 3.1 70B & 44.5 & [39.8, 49.2] & 2.40 \\
Qwen 2.5 72B & 46.8 & [42.0, 51.6] & 2.45 \\
\bottomrule
\end{tabular}
\end{table}

GPT-4o and Claude 3.5 Sonnet are not statistically distinguishable ($p > 0.05$), but both significantly outperform other models ($p < 0.01$).

%==============================================================================
\section{Discussion}
\label{sec:discussion}
%==============================================================================

\paragraph{Personalization Remains Unsolved.} Our results reveal that current frontier models struggle with persona-grounded execution. With memory utilization at 47\% and preference adherence at 61\%, agents are ignoring significant portions of the user context they're provided. This suggests that simply including persona information in the prompt is insufficient---models need architectural or training innovations to properly leverage user-specific context.

\paragraph{Temporal Context is Particularly Hard.} The 14.6-point gap between entity graph usage and session history usage suggests that reasoning over temporal context (past meetings, decisions) is significantly harder than using structured relationships. This may require advances in memory architectures \citep{memgpt2023} or training specifically for temporal reasoning.

\paragraph{Tone Sensitivity Needs Work.} Agents achieving only 58\% on tone appropriateness indicates a failure to adapt communication style to recipient context. In real deployments, an agent sending an insensitive message to a struggling employee could cause significant harm.

\paragraph{Consistency is Critical.} Pass$^3$ scores of 18-32\% mean agents produce the correct behavior in all 3 trials less than a third of the time. For production deployment, such inconsistency is unacceptable.

%==============================================================================
\section{Limitations}
\label{sec:limitations}
%==============================================================================

\paragraph{Persona Scope.} Our 5 personas represent professional knowledge workers. Expanding to other demographics (students, healthcare workers, etc.) would improve generalizability.

\paragraph{Task Scale.} With 20 tasks, our benchmark is smaller than some alternatives. However, the persona-grounded evaluation paradigm means each task is effectively 5 evaluation points (one per persona), yielding 100 distinct evaluations.

\paragraph{LLM Judge Bias.} Using GPT-4o as the tone judge may introduce systematic biases. Future work should validate with human judges or ensemble approaches.

\paragraph{Mock Backends.} Our mock backends simplify real API behavior. Hybrid evaluation with live APIs would improve ecological validity while retaining reproducibility for core metrics.

%==============================================================================
\section{Conclusion}
\label{sec:conclusion}
%==============================================================================

We introduced \klikbench{}, the first benchmark for evaluating AI agents on persona-grounded task execution. Our experiments across 1,500 runs reveal significant gaps: agents utilize only 47\% of relevant memory context, respect tool preferences only 61\% of the time, and achieve appropriate tone in sensitive communications only 58\% of the time. These findings indicate that personalization---adapting agent behavior to individual users---remains an open challenge for AI systems. We release \klikbench{} to support research toward more context-aware, user-adapted AI agents.

%==============================================================================
\section*{Ethics Statement}
%==============================================================================

\klikbench{} evaluates agent personalization using synthetic personas. All entity graph data (people, projects, organizations) is fictional and does not represent real individuals or companies. The ``sensitive recipient'' scenarios (burnout, bereavement) are designed to test appropriate communication, not to exploit vulnerable populations. Mock backends are sandboxed with no external effects. We release under Apache 2.0 to support reproducible research.

%==============================================================================
% References
%==============================================================================

\bibliography{references}

\end{document}
